%%%%%%%%%%  scan idx 467  |  folio 454  %%%%%%%%%%
\textbf{Théorème 4.6.}
Soient \(X\) un schéma quasi-compact ayant un faisceau ample, \(d\) un entier \(>0\) et \(i:Y\to X\) une immersion régulière de codimension \(d\). Alors, pour tout \(k\), on a :
\[
 i_*\bigl((\Fil^kK^\circ(Y))\otimes_{\mathbb Z}\mathbb Q\bigr)
 \subset
 (\Fil^{k+d}K^\circ(X))\otimes_{\mathbb Z}\mathbb Q.
\]

Remarquons tout d'abord que \(i_*(K^\circ(Y))\subset \Fil^1K^\circ(X)\). En effet, \(X-Y\) est dense dans \(X\) : s'il n'en était pas ainsi, \(Y\) contiendrait un ouvert affine, et sur cet ouvert, l'idéal \(J\) définissant \(Y\) serait un nilidéal, donc un idéal localement nilpotent puisqu'il est de type fini. Or ceci est impossible, puisque les \(J^n/J^{n+1}\) sont des \(\mathcal O_Y\)-Modules localement libres de rang \(\ne0\). L'augmentation dans \(K^\circ(X)\) d'un élément de \(i_*(K^\circ(Y))\) est une fonction localement constante sur \(X\); étant nulle sur \(X-Y\), elle est nulle sur \(X\) entier.

Pour tout groupe abélien \(A\), nous poserons \(A_{\mathbb Q}=A\otimes_{\mathbb Z}\mathbb Q\). Soit \(x\in\Fil^kK^\circ(Y)_{\mathbb Q}\) ; on peut écrire \(x\) sous la forme :
\[
  x=\sum a\,\gamma^{i_1}(x_1)\cdots\gamma^{i_n}(x_n),
\]
avec \(\epsilon(x_i)=0\) pour tout \(i\), \(a\in\mathbb Q\), et \(i_1+\cdots+i_n\ge k\). Les \(x_i\) peuvent s'exprimer comme combinaison à coefficients entiers des classes \(E_\alpha\) d'un nombre fini de \(\mathcal O_Y\)-Modules localement libres de type fini \(E_\alpha\). Soit \(R\) la sous-\(\mathbb Q\)-\(\lambda\)-algèbre de \(K^\circ(Y)_{\mathbb Q}\) engendrée par les \(E_\alpha\) et par \(N\). On a alors :
\[
  x\in\Fil^kR.
\]
Par ailleurs, la filtration de \(\lambda\)-anneau de \(R\) est discrète (VI 6.2); soit \(m\) tel que \(\Fil^mR=0\). Enfin, on peut, quitte à décomposer \(X\) en une somme finie d'ouverts disjoints et à faire la démonstration dans chaque ouvert (le \(K^\circ\) commutant aux sommes directes finies), supposer les \(E_\alpha\) de rang constant \(r_\alpha\). L'anneau d'augmentation de \(R\) est alors \(\mathbb Q\).

